\documentclass[10pt]{extarticle}
\usepackage[a4paper,margin=1.5cm]{geometry}
\usepackage[T1]{fontenc}
\usepackage{amsmath,amssymb}
\usepackage{graphicx}
\usepackage[export]{adjustbox}
\usepackage{multicol}
\usepackage{enumitem}
\usepackage{booktabs,array,tabularx}
\setlength{\columnsep}{0.7cm}
\setlength{\parindent}{0pt}
\setlength{\parskip}{3pt}
\pagestyle{empty}
\makeatletter
% figure registry (figdefs.tex is generated by the extract step)
\newcommand{\qfigdef}[4]{\@namedef{qf@file@#1}{#2}\@namedef{qf@w@#1}{#3}\@namedef{qf@wide@#1}{#4}}
\newcommand{\qf@check}[1]{\@ifundefined{qf@file@#1}{\PackageError{score2}{Unknown figure id #1}{}}{}}
\newcommand{\qf@img}[4][]{\includegraphics[width=#2\dimexpr\@nameuse{qf@w@#4}pt\relax,max width=#3,max height=0.45\textheight,#1]{\@nameuse{qf@file@#4}}}
% Full-width questions: finalize wraps every question that contains a wide figure in \qfullbegin...\qfullend,
% so the WHOLE question (stem, figures, options) is set across the page in reading order, then the two columns
% resume. The question's head - its number, stem and first figure (\qfullsplit marks the end of the head) -
% is kept on one page, so a number is never left at a page foot with its figure overleaf; the rest may break
% across pages normally, so no big gaps are left. Without a \qfullsplit, the whole question is kept together
% if it is under 45% of a page. Never write these yourself.
\newif\ifqf@full
\newif\ifqf@boxed
\newsavebox{\qf@box}
\newcommand{\qf@open}{\global\qf@fulltrue\par\global\qf@boxedtrue
  \setbox\qf@box\vbox\bgroup\hsize\textwidth\linewidth\textwidth}
\newcommand{\qf@emit}[1]{\par\egroup\global\qf@boxedfalse
  \ifdim\dimexpr\ht\qf@box+\dp\qf@box\relax<#1\textheight\noindent\usebox{\qf@box}\par\else\unvbox\qf@box\fi}
\newcommand{\qfullsplit}{\ifqf@boxed\qf@emit{0.85}\fi}
\newcommand{\qf@close}{\ifqf@boxed\qf@emit{0.45}\fi\global\qf@fullfalse}
\newcommand{\qfullbegin}{\end{multicols}\qf@open}
\newcommand{\qfullend}{\qf@close\begin{multicols}{2}}
\newcommand{\qfullstart}{\qf@open}
\newcommand{\qfulldone}{\qf@close}
% \qfig[scale]{id}: block figure, centred, directly where it is written (size set by the tool).
% Fallback if a wide figure is not inside a full-width question: break out of the columns in place - no floats.
\newcommand{\qfig}[2][1]{\qf@check{#2}%
  \ifqf@full\qf@blk{#1}{#2}\else
  \ifnum\@nameuse{qf@wide@#2}=1 \expandafter\@firstoftwo\else\expandafter\@secondoftwo\fi
  {\end{multicols}\begin{center}\qf@img{#1}{\textwidth}{#2}\end{center}\begin{multicols}{2}}%
  {\qf@blk{#1}{#2}}\fi}
\newcommand{\qf@blk}[2]{\par\smallskip{\centering\qf@img{#1}{\linewidth}{#2}\par}\smallskip}
% \qfigi[scale]{id}: inline figure (inside option lists / tables / minipages); never breaks columns.
\newcommand{\qfigi}[2][1]{\qf@check{#2}\qf@img[valign=c]{#1}{\linewidth}{#2}}
\makeatother
% structure macros - lint reads these to check completeness and figure placement
\newcommand{\papersec}[1]{\par\bigskip{\centering\Large\bfseries PAPER-#1\par}\medskip}
\newcommand{\partsec}[2]{\par\medskip\noindent{\large\bfseries PART-#1: #2}\par\smallskip}
\newcommand{\sect}[1]{\par\smallskip\noindent\textbf{#1}\par}
\newcommand{\qn}[1]{\par\smallskip\noindent\textbf{#1.}~}
\newcommand{\paraq}[2]{\par\smallskip\noindent\textbf{Paragraph for Questions #1 and #2}\par}
\newenvironment{qoptions}{\begin{itemize}[leftmargin=1.8em,itemsep=1pt,topsep=2pt,label={}]}{\end{itemize}}
\newcommand{\opt}[2]{\item[(#1)] #2}
% figure registry (generated by score2_tools.py finalize)
\qfigdef{p1_phy_s1_q02_1}{Score Advanced - Test 36 (19 May 2024)-img/p1_phy_s1_q02_1.png}{144.4}{0}
\qfigdef{p1_phy_s1_q04_1}{Score Advanced - Test 36 (19 May 2024)-img/p1_phy_s1_q04_1.png}{85.5}{0}
\qfigdef{p1_phy_s2_q08_1}{Score Advanced - Test 36 (19 May 2024)-img/p1_phy_s2_q08_1.png}{88.0}{0}
\qfigdef{p1_phy_s2_q10_1}{Score Advanced - Test 36 (19 May 2024)-img/p1_phy_s2_q10_1.png}{203.3}{0}
\qfigdef{p1_phy_s3_q03_1}{Score Advanced - Test 36 (19 May 2024)-img/p1_phy_s3_q03_1.png}{79.8}{0}
\qfigdef{p1_chem_s1_q01_A}{Score Advanced - Test 36 (19 May 2024)-img/p1_chem_s1_q01_A.png}{47.0}{0}
\qfigdef{p1_chem_s1_q01_B}{Score Advanced - Test 36 (19 May 2024)-img/p1_chem_s1_q01_B.png}{35.5}{0}
\qfigdef{p1_chem_s1_q01_C}{Score Advanced - Test 36 (19 May 2024)-img/p1_chem_s1_q01_C.png}{63.5}{0}
\qfigdef{p1_chem_s1_q01_D}{Score Advanced - Test 36 (19 May 2024)-img/p1_chem_s1_q01_D.png}{73.4}{0}
\qfigdef{p1_chem_s1_q02_A}{Score Advanced - Test 36 (19 May 2024)-img/p1_chem_s1_q02_A.png}{72.8}{0}
\qfigdef{p1_chem_s1_q02_B}{Score Advanced - Test 36 (19 May 2024)-img/p1_chem_s1_q02_B.png}{172.0}{0}
\qfigdef{p1_chem_s1_q02_C}{Score Advanced - Test 36 (19 May 2024)-img/p1_chem_s1_q02_C.png}{119.6}{0}
\qfigdef{p1_chem_s1_q02_D}{Score Advanced - Test 36 (19 May 2024)-img/p1_chem_s1_q02_D.png}{189.2}{0}
\qfigdef{p1_chem_s2_pre_1}{Score Advanced - Test 36 (19 May 2024)-img/p1_chem_s2_pre_1.png}{107.8}{0}
\qfigdef{p1_chem_s2_pre_2}{Score Advanced - Test 36 (19 May 2024)-img/p1_chem_s2_pre_2.png}{102.1}{0}
\qfigdef{p1_chem_s3_q01_1}{Score Advanced - Test 36 (19 May 2024)-img/p1_chem_s3_q01_1.png}{38.5}{0}
\qfigdef{p1_chem_s3_q02_1}{Score Advanced - Test 36 (19 May 2024)-img/p1_chem_s3_q02_1.png}{51.5}{0}
\qfigdef{p1_chem_s3_q02_B}{Score Advanced - Test 36 (19 May 2024)-img/p1_chem_s3_q02_B.png}{56.3}{0}
\qfigdef{p1_chem_s3_q02_C}{Score Advanced - Test 36 (19 May 2024)-img/p1_chem_s3_q02_C.png}{69.2}{0}
\qfigdef{p1_chem_s3_q02_D}{Score Advanced - Test 36 (19 May 2024)-img/p1_chem_s3_q02_D.png}{131.0}{0}
\qfigdef{p1_chem_s3_q02_5}{Score Advanced - Test 36 (19 May 2024)-img/p1_chem_s3_q02_5.png}{205.4}{0}
\qfigdef{p1_chem_s3_q02_6}{Score Advanced - Test 36 (19 May 2024)-img/p1_chem_s3_q02_6.png}{200.6}{0}
\qfigdef{p1_math_s3_q04_1}{Score Advanced - Test 36 (19 May 2024)-img/p1_math_s3_q04_1.png}{108.3}{0}
\qfigdef{p2_phy_s1_q01_1}{Score Advanced - Test 36 (19 May 2024)-img/p2_phy_s1_q01_1.png}{112.2}{0}
\qfigdef{p2_phy_s1_q04_1}{Score Advanced - Test 36 (19 May 2024)-img/p2_phy_s1_q04_1.png}{95.7}{0}
\qfigdef{p2_phy_s2_q05_1}{Score Advanced - Test 36 (19 May 2024)-img/p2_phy_s2_q05_1.png}{85.8}{0}
\qfigdef{p2_phy_s3_para1_2_1}{Score Advanced - Test 36 (19 May 2024)-img/p2_phy_s3_para1_2_1.png}{159.0}{0}
\qfigdef{p2_phy_s3_para3_4_1}{Score Advanced - Test 36 (19 May 2024)-img/p2_phy_s3_para3_4_1.png}{122.6}{0}
\qfigdef{p2_phy_s4_q06_1}{Score Advanced - Test 36 (19 May 2024)-img/p2_phy_s4_q06_1.png}{91.5}{0}
\qfigdef{p2_phy_s4_q07_1}{Score Advanced - Test 36 (19 May 2024)-img/p2_phy_s4_q07_1.png}{76.0}{0}
\qfigdef{p2_phy_s4_q10_1}{Score Advanced - Test 36 (19 May 2024)-img/p2_phy_s4_q10_1.png}{91.5}{0}
\qfigdef{p2_chem_s1_q01_1}{Score Advanced - Test 36 (19 May 2024)-img/p2_chem_s1_q01_1.png}{65.7}{0}
\qfigdef{p2_chem_s1_q01_2}{Score Advanced - Test 36 (19 May 2024)-img/p2_chem_s1_q01_2.png}{42.0}{0}
\qfigdef{p2_chem_s1_q01_3}{Score Advanced - Test 36 (19 May 2024)-img/p2_chem_s1_q01_3.png}{42.2}{0}
\qfigdef{p2_chem_s1_q02_0}{Score Advanced - Test 36 (19 May 2024)-img/p2_chem_s1_q02_0.png}{261.2}{1}
\qfigdef{p2_chem_s1_q02_AL}{Score Advanced - Test 36 (19 May 2024)-img/p2_chem_s1_q02_AL.png}{45.1}{0}
\qfigdef{p2_chem_s1_q02_AP}{Score Advanced - Test 36 (19 May 2024)-img/p2_chem_s1_q02_AP.png}{62.4}{0}
\qfigdef{p2_chem_s1_q02_BL}{Score Advanced - Test 36 (19 May 2024)-img/p2_chem_s1_q02_BL.png}{46.7}{0}
\qfigdef{p2_chem_s1_q02_BP}{Score Advanced - Test 36 (19 May 2024)-img/p2_chem_s1_q02_BP.png}{51.5}{0}
\qfigdef{p2_chem_s1_q02_CL}{Score Advanced - Test 36 (19 May 2024)-img/p2_chem_s1_q02_CL.png}{49.9}{0}
\qfigdef{p2_chem_s1_q02_CP}{Score Advanced - Test 36 (19 May 2024)-img/p2_chem_s1_q02_CP.png}{62.0}{0}
\qfigdef{p2_chem_s1_q02_DL}{Score Advanced - Test 36 (19 May 2024)-img/p2_chem_s1_q02_DL.png}{49.3}{0}
\qfigdef{p2_chem_s1_q02_DP}{Score Advanced - Test 36 (19 May 2024)-img/p2_chem_s1_q02_DP.png}{58.5}{0}
\qfigdef{p2_chem_s1_q03_1}{Score Advanced - Test 36 (19 May 2024)-img/p2_chem_s1_q03_1.png}{138.0}{0}
\qfigdef{p2_chem_s1_q03_A}{Score Advanced - Test 36 (19 May 2024)-img/p2_chem_s1_q03_A.png}{107.8}{0}
\qfigdef{p2_chem_s1_q03_B}{Score Advanced - Test 36 (19 May 2024)-img/p2_chem_s1_q03_B.png}{87.6}{0}
\qfigdef{p2_chem_s1_q03_C}{Score Advanced - Test 36 (19 May 2024)-img/p2_chem_s1_q03_C.png}{75.3}{0}
\qfigdef{p2_chem_s1_q03_D}{Score Advanced - Test 36 (19 May 2024)-img/p2_chem_s1_q03_D.png}{75.0}{0}
\qfigdef{p2_chem_s1_q04_A}{Score Advanced - Test 36 (19 May 2024)-img/p2_chem_s1_q04_A.png}{138.7}{0}
\qfigdef{p2_chem_s1_q04_B}{Score Advanced - Test 36 (19 May 2024)-img/p2_chem_s1_q04_B.png}{141.1}{0}
\qfigdef{p2_chem_s1_q04_C}{Score Advanced - Test 36 (19 May 2024)-img/p2_chem_s1_q04_C.png}{140.4}{0}
\qfigdef{p2_chem_s1_q04_D}{Score Advanced - Test 36 (19 May 2024)-img/p2_chem_s1_q04_D.png}{140.2}{0}
\qfigdef{p2_chem_s3_para1_2_1}{Score Advanced - Test 36 (19 May 2024)-img/p2_chem_s3_para1_2_1.png}{201.3}{0}
\qfigdef{p2_chem_s3_para3_4_1}{Score Advanced - Test 36 (19 May 2024)-img/p2_chem_s3_para3_4_1.png}{243.2}{0}
\qfigdef{p2_chem_s4_q09_1}{Score Advanced - Test 36 (19 May 2024)-img/p2_chem_s4_q09_1.png}{57.6}{0}
\raggedcolumns
\makeatletter\@beginparpenalty=10000\makeatother
\newcolumntype{L}[1]{>{\raggedright\arraybackslash}p{#1}}
\begin{document}
\begin{center}{\Large\bfseries Score Advanced --- Test 36 (19 May 2024)}\end{center}
\begin{multicols}{2}
\raggedcolumns
\papersec{1}
\partsec{1}{PHYSICS}
\sect{SECTION-I(i)}
\qn{1} A variable capacitor with initial capacitance C$_0$ is charged to a voltage V$_0$ and is short circuited through a resistance R. If current remains constant during discharging (an external agent is pulling the plates apart) then choose correct option(s).
\begin{qoptions}
\opt{A}{Variation of capacitance with time is given by C = C$_0\left(1-\dfrac{t}{RC_0}\right)$.}
\opt{B}{Time taken to discharge the capacitor is given by t = RC$_0$.}
\opt{C}{Total heat produced in the resistor is C$_0$V$_0^2$.}
\opt{D}{Initial energy stored in the capacitor is $\dfrac12$C$_0$V$_0^2$.}
\end{qoptions}
\qn{2} A cylindrical body of radius R is fixed in a horizontal position. A piece of ideal thread of length $\ell$ was laid on its slippery surface. On one end of thread a point like body of mass m is attached, while to the other end another point like body of mass 2m is attached . Choose correct statement(s). The system is in equilibrium.\par
\qfig{p1_phy_s1_q02_1}
\begin{qoptions}
\opt{A}{2 sin$\alpha$ = sin$\beta$}
\opt{B}{2 sin$\beta$ = sin$\alpha$}
\opt{C}{Maximum possible length of string is $\ell_{\max}$ = 2$\pi$R/3.}
\opt{D}{Maximum possible length of string is $\ell_{\max}$ = 2R.}
\end{qoptions}
\qn{3} By properly combining two prisms made of different materials, it is possible to :-
\begin{qoptions}
\opt{A}{have dispersion without average deviation.}
\opt{B}{have deviation without dispersion.}
\opt{C}{have both dispersion and average deviation.}
\opt{D}{have neither dispersion nor average deviation.}
\end{qoptions}
\qn{4} A sample of ideal gas of some mass is taken through the cyclic process shown in the figure. The temperature of the gas at state A is T$_\mathrm A$ = 200 K, and at state C is T$_\mathrm C$ = 1800 K. Choose the correct option(s) :-\par
\qfig{p1_phy_s1_q04_1}
\begin{qoptions}
\opt{A}{$\dfrac{V_A}{V_C}=\dfrac13$}
\opt{B}{$\dfrac{V_A}{V_B}=\dfrac19$}
\opt{C}{$\dfrac{P_AV_A}{P_CV_C}=\dfrac13$}
\opt{D}{$\dfrac{P_AV_A}{P_CV_C}=\dfrac19$}
\end{qoptions}
\qn{5} X-ray from a tube with a target A of atomic number Z shows strong K lines for target A and weak K lines for impurities. The wavelength of K$_\alpha$ lines is $\lambda_z$ for target A and $\lambda_1$ and $\lambda_2$ for two impurities. If $\dfrac{\lambda_z}{\lambda_1}$ = 4 and $\dfrac{\lambda_z}{\lambda_2}=\dfrac14$. Screening constant of K$_\alpha$ lines to be unity. Select the correct statement(s).
\begin{qoptions}
\opt{A}{The atomic number of first impurity is 2Z $-$ 1.}
\opt{B}{The atomic number of first impurity is 2Z + 1.}
\opt{C}{The atomic number of second impurity is $\dfrac{(Z+1)}2$.}
\opt{D}{The atomic number of second impurity is $\dfrac Z2$ + 1.}
\end{qoptions}
\qn{6} A man is sitting in a boat floating in water. There are heavy stones placed in the boat.
\begin{qoptions}
\opt{A}{When the man throws the stones in water, the level of boat goes down.}
\opt{B}{When the man throws the stones in water, the level of boat rises up.}
\opt{C}{When the man drinks some water, the level of boat goes down}
\opt{D}{When the man drinks some water, the level of boat remains unchanged.}
\end{qoptions}
\sect{SECTION-I(ii)}
\qn{7} A block of mass 1.0 kg is pulled up on a rough slope of inclination 37$^\circ$ with the horizontal. The coefficient of kinetic friction between the block and the slope is 0.25. The block moves with acceleration 2.0 m/s$^2$. After 5.0 s from the instant the block starts, which of the following statements is/are true ? (Power of any force is given by P = $\vec F\cdot\vec V$)\par
{\footnotesize\setlength{\tabcolsep}{1.5pt}\noindent\begin{tabular}{|c|L{3.6cm}|c|L{2.0cm}|}
\hline
 & \multicolumn{1}{>{\centering\arraybackslash}p{3.6cm}|}{List-I} & & \multicolumn{1}{>{\centering\arraybackslash}p{2.0cm}|}{List-II} \\ \hline
(I) & Power delivered by the pulling force & (P) & 20 W \\ \hline
(II) & Power delivered by the weight of the block & (Q) & $-$20 W \\ \hline
(III) & Power delivered by the friction between the slope and the block & (R) & $-$60 W \\ \hline
(IV) & Power delivered by the contact force between the slope and the block & (S) & 100 W \\ \hline
 & & (T) & 80 W \\ \hline
\end{tabular}}\par
\begin{qoptions}
\opt{A}{I $\to$ P;II $\to$ R;III $\to$ Q;IV $\to$ S}
\opt{B}{I $\to$ S;II $\to$ R;III $\to$ Q;IV $\to$ Q}
\opt{C}{I $\to$ P;II $\to$ Q;III $\to$ R;IV $\to$ S}
\opt{D}{I $\to$ Q;II $\to$ S;III $\to$ T;IV $\to$ P}
\end{qoptions}
\qn{8} A disc of mass m and radius R is performing pure rolling on the inclined surface of a wedge of mass M. Wedge is placed on a rough horizontal surface. If wedge does not move on the horizontal surface then match the followings.\par
\qfig{p1_phy_s2_q08_1}
{\footnotesize\setlength{\tabcolsep}{1.5pt}\noindent\begin{tabular}{|c|L{2.0cm}|c|L{4.0cm}|}
\hline
 & \multicolumn{1}{>{\centering\arraybackslash}p{2.0cm}|}{List-I} & & \multicolumn{1}{>{\centering\arraybackslash}p{4.0cm}|}{List-II} \\ \hline
(I) & Frictional force acting on disc & (P) & $\dfrac{mg\sin\theta}{3}$ \\ \hline
(II) & Vertical component of net frictional force acting on wedge & (Q) & $\dfrac{mg\sin^2\theta}{3}$ \\ \hline
(III) & Normal reaction acting on disc & (R) & mg$\cos\theta$ \\ \hline
(IV) & Normal reaction acting on wedge by the ground & (S) & $\left(m\cos^2\theta+\dfrac m3\sin^2\theta+M\right)g$ \\ \hline
 & & (T) & $\left(m\sin^2\theta+\dfrac m3\cos^2\theta+M\right)g$ \\ \hline
\end{tabular}}\par
\begin{qoptions}
\opt{A}{I $\to$ P;II $\to$ Q;III $\to$ S;IV $\to$ T}
\opt{B}{I $\to$ P;II $\to$ Q;III $\to$ T;IV $\to$ S}
\opt{C}{I $\to$ P;II $\to$ Q;III $\to$ R;IV $\to$ S}
\opt{D}{I $\to$ S;II $\to$ T;III $\to$ P;IV $\to$ Q}
\end{qoptions}
\qn{9} Match the following lists.\par
{\footnotesize\setlength{\tabcolsep}{1.5pt}\noindent\begin{tabular}{|c|L{3.0cm}|c|L{3.0cm}|}
\hline
 & \multicolumn{1}{>{\centering\arraybackslash}p{3.0cm}|}{List$-$I} & & \multicolumn{1}{>{\centering\arraybackslash}p{3.0cm}|}{List$-$II} \\ \hline
(I) & Nearby objects appear moving past faster than distant objects. & (P) & Human eye sense the direction from which light enters it. \\ \hline
(II) & Extremely distant objects appear moving among with us & (Q) & It cannot resolve two points if they subtend an angle less than one minute. \\ \hline
(III) & Many people believe that they can read better in bright light than in dim light & (R) & Persistence of human eye is 0.1 s \\ \hline
(IV) & The colors of the surface of bodies that can be impregnated with water appear richer after moistening the surface. & (S) & Contraction of the pupil in bright light reduces spherical aberration. \\ \hline
 & & (T) & Due to moistening, surface irregularities are covered with water film, which reduces diffuse reflection of light \\ \hline
\end{tabular}}\par
\begin{qoptions}
\opt{A}{I $\to$ P;II $\to$ Q;III $\to$ R;IV $\to$ S}
\opt{B}{I $\to$ P;II $\to$ T;III $\to$ R;IV $\to$ Q}
\opt{C}{I $\to$ P;II $\to$ S;III $\to$ R;IV $\to$ P}
\opt{D}{I $\to$ P;II $\to$ P;III $\to$ S;IV $\to$ T}
\end{qoptions}
\qn{10} \qfig{p1_phy_s2_q10_1}
The right limb is long enough and its radius is twice that of left limb. Water is slowly poured into the right limb. Considering water wets completely and the left limb is fully filled. List I corresponds to height of liquid in right limb and list II corresponds to radius of meniscus (R) and its shape in left limb.\par
{\footnotesize\setlength{\tabcolsep}{1.5pt}\noindent\begin{tabular}{|c|L{2.4cm}|c|L{2.4cm}|}
\hline
 & \multicolumn{1}{>{\centering\arraybackslash}p{2.4cm}|}{List-I} & & \multicolumn{1}{>{\centering\arraybackslash}p{2.4cm}|}{List-II} \\ \hline
(P) & h = 1.1 cm & (1) & Concave upward \\ \hline
(Q) & h = 2 cm & (2) & Flat \\ \hline
(R) & h = 2.54 cm & (3) & Convex upward \\ \hline
(S) & h = 6.86 cm & (4) & R = 1mm \\ \hline
 & & (5) & R = 0.5 mm \\ \hline
\end{tabular}}\par
\begin{qoptions}
\opt{A}{P $\to$ 1,4;Q $\to$ 1;R $\to$ 2;S $\to$ 3,4}
\opt{B}{P $\to$ 1;Q $\to$ 1;R $\to$ 2;S $\to$ 3,4}
\opt{C}{P $\to$ 1,5;Q $\to$ 1;R $\to$ 1,4;S $\to$ 3,5}
\opt{D}{P $\to$ 1,5;Q $\to$ 1;R $\to$ 1,5;S $\to$ 3}
\end{qoptions}
\sect{SECTION-II}
\qn{1} In a long cylindrical metal wire of radius r and of resistivity $\rho$, a current of strength I flows in uniform distribution. The wire has constant surface temperature T$_0$. The temperature of the wire on its axis of symmetry is given by T = T$_\mathrm o$ + $\dfrac{I^2\rho}{\pi^2r^2K}\times\lambda$. Here K is thermal conductivity of wire. Find $\lambda$.
\qn{2} A certain amount of water is heated from 20$^\circ$C to 40$^\circ$C using an 800 Watt immersion heater. We expect to heat water in 210 seconds, but instead it takes 230 seconds. Determine the heat capacity of vessel in S.I. units. Ignore other heat losses.
\qn{3} A mass M is divided in two parts m and M$-$m. Now two parts are connected by an ideal string passing over a pulley as shown in diagram. Find value of m/M such that tension in string is maximum.\par
\qfig{p1_phy_s3_q03_1}
\qn{4} A cell of emf 1.5 V and internal resistance 0.5 $\Omega$ is connected to a nonlinear conductor in which current varies with voltage as V$^2$ = I. Calculate the current drawn from cell (in ampere).
\qn{5} A coil with inductance L and resistance R is connected to an alternating source. The capacity of a capacitor that can be connected in series later with the coil such that active power of circuit does not change is given by\par C = $\dfrac{n}{\omega^2L}$. Find value of n.
\qn{6} Consider a spherical homogeneous cloud of mass M made by an explosion. Due to energy received from the explosion, it is expanding and the expansion is spherically symmetric. At an instant its radius is R$_0$ and the particles on the surface are moving away from its center with velocity V$_0$. Its radius when expansion ceases is given by\par $\dfrac{nGMR_0}{2GM-R_0V_0^2}$. Find n.
\qn{7} The electrostatic flux through a rectangular area of $\dfrac7{188}$m$^2$ is lying in a plane 2x + 3y + 6z = 10 due to an electric field (8$\hat i$ + 6$\hat j$ + 10$\hat k$) V/m will be in (V-m)
\qn{8} A parallel beam of light of all wavelength greater than 3000 \AA\ falls on a double slit in a Young's double slit experiment. It is observed that the wavelengths 3600\AA\ and 6000\AA\ are absent at a distance of 31.5 mm from the position of the central maximum, and the orders of the minima at this point for the two wavelengths differ by 7. If the distance between the slits and the screen is 1m, the separation between the two slits (in mm) is :
\raggedcolumns
\papersec{1}
\partsec{2}{CHEMISTRY}
\sect{SECTION-I(i)}
\qn{1} A compound `X' evolves CO$_2$ when treated with NaHCO$_3$ but gives negative test with tollen's reagent and fehling's solution. On treatment with CHCl$_3$, KOH, a foul smelling organic compound is formed. If `X' do not give positive iodoform test, which of the following satisfy the conditions to be `X'.
\begin{qoptions}
\opt{A}{\qfigi{p1_chem_s1_q01_A}}
\opt{B}{\qfigi{p1_chem_s1_q01_B}}
\opt{C}{\qfigi{p1_chem_s1_q01_C}}
\opt{D}{\qfigi{p1_chem_s1_q01_D}}
\end{qoptions}
\qn{2} Which of the following method(s) is/are not used to prepare p-bromo aniline as major product :
\begin{qoptions}
\opt{A}{\qfigi{p1_chem_s1_q02_A}}
\opt{B}{\qfigi{p1_chem_s1_q02_B}}
\opt{C}{\qfigi{p1_chem_s1_q02_C}}
\opt{D}{\qfigi{p1_chem_s1_q02_D}}
\end{qoptions}
\qn{3} In which of the following process(es) electrons are removed from different type of orbitals?
\begin{qoptions}
\opt{A}{O$_2^-\to$ O$_2^+$}
\opt{B}{Cu $\to$ Cu$^{2+}$}
\opt{C}{N$_2^+\to$ N$_2^-$}
\opt{D}{NO $\to$ NO$^{2+}$}
\end{qoptions}
\qn{4} Which of the following statements are incorrect from following options?
\begin{qoptions}
\opt{A}{Zone refining is carried out in presence of noble gas environment.}
\opt{B}{Fe is the 2$^{\mathrm{nd}}$ most abundant metal in the earth crust.}
\opt{C}{Sapphire consists of Al$_2$O$_3$ and Co (as impurity).}
\opt{D}{Formula of Bauxite : AlO$_{3-2x}$(OH)$_x$ , where 0 $<$ x $<$ 1.}
\end{qoptions}
\qn{5} Which of following is/are incorrect statement(s) for a sample of ideal gas ?
\begin{qoptions}
\opt{A}{Temperature of gas remains constant during its adiabatic free expansion}
\opt{B}{During compression of the gas at constant pressure, temperature of gas increases.}
\opt{C}{Molar heat capacity of the gas is zero during an isothermal process}
\opt{D}{At 300 K temperature molar internal energy change of SO$_2$(g) is greater than molar internal energy change of CO$_2$(g) when temperature of both gases is increased by 20 K.}
\end{qoptions}
\qn{6} Select the correct statement(s) about FCC(ABCABC...) structure,\par
[Given : a = Edge length of FCC unit cell]
\begin{qoptions}
\opt{A}{Distance between nearest octahedral void and tetrahedral void is $\dfrac{\sqrt3a}4$}
\opt{B}{Distance between two nearest octahedral voids is $\dfrac a{\sqrt2}$}
\opt{C}{Distance between two nearest tetrahedral voids is $\dfrac{\sqrt3a}2$}
\opt{D}{Distance between close-packed layers A and B is 2r$\sqrt{\dfrac23}$ [r = radius of atom]}
\end{qoptions}
\sect{SECTION-I(ii)}
\qn{7} {\footnotesize\setlength{\tabcolsep}{1.5pt}\noindent\begin{tabular}{|c|L{3.6cm}|c|L{2.6cm}|}
\hline
\multicolumn{2}{|>{\centering\arraybackslash}p{3.9cm}|}{List-I (Compound)} & \multicolumn{2}{>{\centering\arraybackslash}p{3.0cm}|}{List-II (Can be distinguish by)} \\ \hline
(P) & CH$_3-$CHO \& Ph$-$CHO & (1) & Br$_2$ / H$_2$O test \\ \hline
(Q) & CH$_3-$CH$_3$ \& HC$\equiv$CH & (2) & Ammonical Aq. NO$_3^-$ \\ \hline
(R) & \qfigi{p1_chem_s2_pre_1} & (3) & Iodoform test \\ \hline
(S) & \qfigi{p1_chem_s2_pre_2} & (4) & NaHSO$_4$ test \\ \hline
 & & (5) & Fehling test \\ \hline
\end{tabular}}\par
\begin{qoptions}
\opt{A}{P $\to$ 3,5;Q $\to$ 1,2;R $\to$ 1,2,3,4;S $\to$ 2,5}
\opt{B}{P $\to$ 1,2;Q $\to$ 2,3;R $\to$ 3,4;S $\to$ 4,5}
\opt{C}{P $\to$ 3,4,5;Q $\to$ 2,3,4;R $\to$ 4,5;S $\to$ 1,2,3}
\opt{D}{P $\to$ 3,4;Q $\to$ 2,4;R $\to$ 1,2;S $\to$ 4,5}
\end{qoptions}
\qn{8} List-I contains co-ordination compounds and List-II contains their properties:\par
{\footnotesize\setlength{\tabcolsep}{1.5pt}\noindent\begin{tabular}{|c|L{2.4cm}|c|L{3.6cm}|}
\hline
 & \multicolumn{1}{>{\centering\arraybackslash}p{2.4cm}|}{List-I} & & \multicolumn{1}{>{\centering\arraybackslash}p{3.6cm}|}{List-II} \\ \hline
(P) & [NiCl$_4$]$^{2-}$ & (1) & Diamagnetic species \\ \hline
(Q) & [Cr(H$_2$O)$_3$F$_3$] & (2) & Low spin octahedral complex \\ \hline
(R) & [PtCl$_4$]$^{2-}$ & (3) & Contains unpaired electrons \\ \hline
(S) & [Co(H$_2$O)$_6$]$^{3+}$ & (4) & d$^8$ configuration of metal ion \\ \hline
 & & (5) & Exhibits geometrical isomerism \\ \hline
\end{tabular}}\par
[Given atomic no. Ni = 28, Cr = 24, Pt = 78, Co = 27]\par
Match each species in List-I with their properties in List-II and choose the correct option:
\begin{qoptions}
\opt{A}{P $\to$ 1,2;Q $\to$ 3,4;R $\to$ 3,5;S $\to$ 1,4}
\opt{B}{P $\to$ 3,4;Q $\to$ 2,5;R $\to$ 1,4;S $\to$ 1,2}
\opt{C}{P $\to$ 4,5;Q $\to$ 1,2;R $\to$ 2,3;S $\to$ 3,5}
\opt{D}{P $\to$ 3,4;Q $\to$ 3,5;R $\to$ 1,4;S $\to$ 1,2}
\end{qoptions}
\qn{9} Boyle's temperature (T$_\mathrm B$): T$_\mathrm B$ of CO$_2$ = 350$^\circ$C, T$_\mathrm B$ of CH$_4$ = 60$^\circ$C and T$_\mathrm C$ of H$_2$ = $-$240$^\circ$C\par
{\footnotesize\setlength{\tabcolsep}{1.5pt}\noindent\begin{tabular}{|c|L{3.0cm}|c|L{3.0cm}|}
\hline
 & \multicolumn{1}{>{\centering\arraybackslash}p{3.0cm}|}{List-I} & & \multicolumn{1}{>{\centering\arraybackslash}p{3.0cm}|}{List-II} \\ \hline
(P) & CO$_2$ gas in low pressure region \& temperature 573 K. & (1) & PV = nRT \\ \hline
(Q) & CO$_2$ gas in very high pressure region \& temperature 623 K & (2) & Gas is more compressible with respect to an ideal gas. \\ \hline
(R) & CH$_4$ gas at 333 K \& in low pressure region. & (3) & Gas cannot be liquefied at any pressure at given temperature. \\ \hline
(S) & H$_2$ gas in low pressure region \& temperature 20 K. & (4) & Molar volume of gas is lesser than that of ideal gas at same pressure and temperature. \\ \hline
 & & (5) & Repulsive forces dominate \\ \hline
\end{tabular}}\par
\begin{qoptions}
\opt{A}{P $\to$ 2,3;Q $\to$ 1,2;R $\to$ 1,2;S $\to$ 2,4}
\opt{B}{P $\to$ 2,3,4,5;Q $\to$ 3,5;R $\to$ 1,3;S $\to$ 2,4}
\opt{C}{P $\to$ 2,3,4;Q $\to$ 3,5;R $\to$ 1,3;S $\to$ 1,2,4}
\opt{D}{P $\to$ 2,3,4;Q $\to$ 3,5;R $\to$ 1,3;S $\to$ 2,4}
\end{qoptions}
\qn{10} Match the following\par
{\footnotesize\setlength{\tabcolsep}{1.5pt}\noindent\begin{tabular}{|c|L{3.0cm}|c|L{3.0cm}|}
\hline
 & \multicolumn{1}{>{\centering\arraybackslash}p{3.0cm}|}{List - I} & & \multicolumn{1}{>{\centering\arraybackslash}p{3.0cm}|}{List - II} \\ \hline
(P) & A saturated aqueous solution of Mg(OH)$_2$ [K$_\mathrm{sp}$ = 4$\times$10$^{-12}$] & (1) & pH = 7 \\ \hline
(Q) & 300 ml of 0.1 M BOH +100 ml. of 0.1M HCl [K$_\mathrm b$(BOH) = 10$^{-4}$] & (2) & pH = $\dfrac{\mathrm{pKa_2+pKa_3}}2$ \\ \hline
(R) & 300 ml. of 0.1 M BOH+300 ml. of 0.1 M HA [K$_\mathrm a$ =K$_\mathrm b$ = 10$^{-4}$] & (3) & pH= 10 + log 2 \\ \hline
(S) & 0.01 M aqueous solution Na$_2$HPO$_4$ [Given that: pKa$_1$ (H$_3$PO$_4$) = 4.3; pKa$_2$ (H$_3$PO$_4$) = 8.102; pKa$_3$ (H$_3$PO$_4$) = 12.5; & (4) & pH = $\dfrac12$pK$_\mathrm w$ + $\dfrac12$pK$_\mathrm a$ $-$ $\dfrac12$pK$_\mathrm b$ \\ \hline
 & & (5) & Buffer solution \\ \hline
\end{tabular}}\par
\begin{qoptions}
\opt{A}{P $\to$ 3;Q $\to$ 3;R $\to$ 2;S $\to$ 3}
\opt{B}{P $\to$ 1;Q $\to$ 3;R $\to$ 1;S $\to$ 4}
\opt{C}{P $\to$ 3;Q $\to$ 1;R $\to$ 4;S $\to$ 3}
\opt{D}{P $\to$ 3;Q $\to$ 3;R $\to$ 5;S $\to$ 2}
\end{qoptions}
\sect{SECTION-II}
\qn{1} Total number of fractions obtained on fractional distillation of all monochloro derivatives product(s)\par
\qfig{p1_chem_s3_q01_1}
\qn{2} \noindent(A)~\qfigi{p1_chem_s3_q02_1}\hfill(B)~\qfigi{p1_chem_s3_q02_B}\par
\noindent(C)~\qfigi{p1_chem_s3_q02_C}\par
\noindent(D)~\qfigi{p1_chem_s3_q02_D}\par
\qfig{p1_chem_s3_q02_5}
\qfig{p1_chem_s3_q02_6}
Calculate the number of reducing sugar = x\par
Calculate the number of sugar in which molecular weight increases by 210, whenever it undergo acetylation reaction with acetic anhydride = y\par
So x $-$ y =
\qn{3} Organic compound [M] has the following data :\par
(i) Molecular weight is 149, it contains only 3 different elements.\par
(ii) It can rotate plane polarized light (chiral carbon).\par
(iii) [M] on oxidation with hot alkaline KMnO$_4$ produces benzoic acid.\par
(iv) [M] can be prepared by treating carbonyl compound with methyl amine followed by NaBH$_4$.\par
Write the number of possible structures for [M] (excluding stereo isomers).
\qn{4} A mixture of Potash alum, Chrome alum and Mohr's salt is taken in water and H$_2$O$_2$ is added in strongly alkaline medium and finally filtered.\par
If atomic mass/molecular mass of cation in residue is ``a'' and the number of type of ions from the salt mixture passed into the filtrate is `b'. Then find the value of a/b.\par
(Atomic mass : Cr = 52 , Fe = 56 , S = 32 , O = 16 , N = 14 , C = 12 , B = 11 , H = 1 , K$^+$ = 39 , Na$^+$ = 23 , Ni = 58 , Cu = 29 , Zn = 65 , Cl = 35.5 , Br = 80 , F = 19)
\qn{5} (a) Fullerene has no dangling bonds.\par
(b) Ni + excess NH$_3$ solution $\longrightarrow$ deep blue solution\par
(c) PH$_3$ can be absorbed by HgCl$_2$(aq.)\par
(d) Bleaching powder available in market is completely soluble in water.\par
(e) Na can be safely stored in ethanol.\par
(f) Na$_2$S$_2$O$_3$ + Cl$_2$-water (excess) produces NaHSO$_4$ and HCl.\par
(g) Higher concentration of SO$_2$ damages the leaves of plant and retard the rate of photosynthesis.\par
If x = number of correct statements.\par
y = number of incorrect statements.\par
then find the value of x/y.
\qn{6} Let `A' and `B' are two completely miscible liquids. P$^\circ_\mathrm A$ = 90 torr and P$^\circ_\mathrm B$ = 40 torr. In a flask 8 moles of `A' is mixed with 4 moles of `B'. However as soon as B is added `A' starts dissociating as two different liquids `C' and `D' as A $\longrightarrow$ C + D.\par
Both `C' and `D' are also completely miscible with each other and with `A' and `B'. Given P$^\circ_\mathrm C$ = 50 torr and P$^\circ_\mathrm D$ = 40 torr.\par
Assume dissociation of `A' follows 1$^{\mathrm{st}}$ order kinetics with half life equal to 50 min. After 100 min 2 moles of a non-volatile non-electrolyte solute `S' is added in solution which immediately arrests the dissociation of A completely.\par
Calculate final vapour pressure of the solution in torr.
\qn{7} Fixed amount of an ideal gas contained in a sealed rigid vessel of volume 24.6 litres at 1. 0 bar is heated reversibly from 27$^\circ$C to 127$^\circ$C. The entropy of gas is S = 10 + 10$^{-2}$ T (J/K). Then what is the enthalpy change (in J) during the process.
\qn{8} NO$_3^-\to$ NO$_2$(acid medium), E$^\circ$=0.790V\par
NO$_3^-\to$ NH$_3$OH$^+$(acid medium) E$^\circ$=0.731V.\par
At what pH. the above two electrodes will have same electrode potential value?\par
Assume that activity of all other species to be unity except [H$^+$] ions ? $\left(\dfrac{2.303\,RT}{F}=0.059\right)$
\raggedcolumns
\papersec{1}
\partsec{3}{MATHEMATICS}
\sect{SECTION-I(i)}
\qn{1} Let the equation ax$^2-$bx + c = 0 has 2 distinct roots in the interval (0,1) where a,b,c $\in$ N. If $\lambda<\log_5$(abc) for all choices of natural numbers a,b,c then non-negative integral values of $\lambda$ can be :
\begin{qoptions}
\opt{A}{0}
\opt{B}{1}
\opt{C}{2}
\opt{D}{3}
\end{qoptions}
\qn{2} Event `A' is independent of event B,B $\cup$ C and B $\cap$ C. If P(A) = $\dfrac12$, P(B) = $\dfrac13$ and P(C) = $\dfrac14$. Then :\par
(Where P(x) denotes probability of event x)
\begin{qoptions}
\opt{A}{P$\left(\dfrac AC\right)=\dfrac12$}
\opt{B}{P$\left(\dfrac{\overline B\cup\overline C}{A}\right)=\dfrac{11}{12}$ (where B and C independent events)}
\opt{C}{P$\left(\dfrac{\overline A}{\overline{B\cap C}}\right)=\dfrac12$}
\opt{D}{A and C are not independent events}
\end{qoptions}
\qn{3} An ellipse with eccentricity $\dfrac12$ passes through P(3,4) whose nearer focus is S(0,0) and equation of tangent at P on ellipse is 3x + 4y $-$ 25 = 0. If a chord through S parallel to tangent at P intersects the ellipse at A and B then :
\begin{qoptions}
\opt{A}{length of AB is 15}
\opt{B}{length of latus rectum of ellipse is 15}
\opt{C}{focal length of ellipse is 10}
\opt{D}{centre of ellipse is ($-$3,$-$4)}
\end{qoptions}
\qn{4} Let $<$T$_n>$ be a sequence such that T$_n^3$ + 2T$_n$ = T$_{n+1}$ $\forall$ n $\in$ N, and T$_1$ = 1, then :
\begin{qoptions}
\opt{A}{$\displaystyle\sum_{n=1}^{100}\left(T_n^3+T_n+1\right)$ = T$_{101}$ + 99}
\opt{B}{$\displaystyle\sum_{n=1}^{100}\left(T_n^3+T_n+1\right)$ = T$_{101}$ + 100}
\opt{C}{$\displaystyle\prod_{n=1}^{100}\left(T_n^2+2\right)$ = T$_{100}$}
\opt{D}{$\displaystyle\prod_{n=1}^{100}\left(\dfrac{T_{n+1}}{T_n}-T_n^2\right)$ = 2$^{100}$}
\end{qoptions}
\qn{5} Let f : (0, $\infty$) $\longrightarrow$ R be a differentiable function satisfying the equation\par
f(xy) = e$^{xy-x-y}$(e$^y$f(x) + e$^x$f(y)) $\forall$ x, y $>$ 0. If f$\,'$(1) = e, then\par
(Note : where `e' is nepier's constant)
\begin{qoptions}
\opt{A}{$\displaystyle\lim_{x\to e}\left[\dfrac{f(x)-e^x}{x-e}\right]$ = e$^{(e-1)}$}
\opt{B}{number of roots of the equation f(x) = xe$^x$ in (0,$\infty$) is 2.}
\opt{C}{$\displaystyle\int_1^e$ f(x) dx $<$ e$^e$ (e $-$ 1)}
\opt{D}{f(x) is strictly increasing function in (0,$\infty$).}
\end{qoptions}
\qn{6} If three planes P$_1\equiv$ 2x + y + z $-$ 1 = 0, P$_2\equiv$ x $-$ y + z $-$ 2 = 0 and P$_3\equiv\alpha$x $-$ y + 3z $-$ 5 = 0 intersects each other at point P on XOY plane and at point Q on YOZ plane, where O is the origin then identify the correct statement(s) ?
\begin{qoptions}
\opt{A}{The value of $\alpha$ is 4.}
\opt{B}{Straight line perpendicular to plane P$_3$ and passing through P is $\dfrac{x-1}{4}=\dfrac{y+1}{-1}=\dfrac z3$.}
\opt{C}{The length of projection of $\overrightarrow{PQ}$ on x-axis is 1.}
\opt{D}{Centroid of the triangle OPQ is $\left(\dfrac13,\dfrac{-1}2,\dfrac12\right)$}
\end{qoptions}
\sect{SECTION-I(ii)}
\qn{7} Match the following List-I with List-II\par
{\footnotesize\setlength{\tabcolsep}{1.5pt}\noindent\begin{tabular}{|c|L{3.7cm}|c|c|}
\hline
 & \multicolumn{1}{>{\centering\arraybackslash}p{3.7cm}|}{List-I} & & List-II \\ \hline
(I) & The number of all triples (a,b,c) such that all three term a,b,c are in harmonic progression in which a = 20 and b divides c (where a,b,c are strictly increasing positive integers ) & (P) & 66 \\ \hline
(II) & The number of solutions of m + n + p = 10, (m,n,p are non-negative integers) & (Q) & 20 \\ \hline
(III) & The number of ordered pair (a,b) satisfying the equation ab + b $-$ a + 1 = 0 are where a,b $\in$ I & (R) & 4 \\ \hline
(IV) & The number of ordered pairs (x,y) of real numbers such that (x + iy)$^{20}$ = x $-$ iy is (where i = $\sqrt{-1}$) & (S) & 5 \\ \hline
 & & (T) & 22 \\ \hline
\end{tabular}}\par
\begin{qoptions}
\opt{A}{I $\to$ S;II $\to$ P;III $\to$ R;IV $\to$ T}
\opt{B}{I $\to$ T;II $\to$ Q;III $\to$ S;IV $\to$ P}
\opt{C}{I $\to$ Q;II $\to$ P;III $\to$ R;IV $\to$ T}
\opt{D}{I $\to$ R;II $\to$ S;III $\to$ Q;IV $\to$ P}
\end{qoptions}
\qn{8} Match the following List-I with List-II\par
{\footnotesize\setlength{\tabcolsep}{1.5pt}\noindent\begin{tabular}{|c|L{4.6cm}|c|c|}
\hline
 & \multicolumn{1}{>{\centering\arraybackslash}p{4.6cm}|}{List-I} & & List-II \\ \hline
(I) & Let f(x) = x$^4$ + ax$^3$ + bx$^2$ + cx + d (where a,b,c,d are real coefficient) and f(x) = 0 has real roots, If $|$f(i)$|$ = 1 and a = b = c = d = k then k can be (where i = $\sqrt{-1}$) & (P) & 0 \\ \hline
(II) & If x$^5$ = 1 (x $\neq$ 1) then $\dfrac{x}{1+x^2}+\dfrac{x^2}{1+x^4}+\dfrac{x^3}{1+x}+\dfrac{x^4}{1+x^3}$ equals & (Q) & 1 \\ \hline
(III) & Let f(x) = $\ell$n$\left[\cos|x|+\dfrac12\right]$ (where [.] denotes the greatest integer function), then $\displaystyle\int_{x_1}^{x_2}\left(\lim_{n\to\infty}\left(\dfrac{(f(x))^n}{x^2+\tan^2x}\right)\right)$dx is (where x$_1$x$_2\in\left[-\dfrac\pi6,\dfrac\pi6\right]$) & (R) & 2 \\ \hline
(IV) & Let ABCD be a cyclic quadrilateral inscribed in a circle of radius 2, such that BD = 2$\sqrt3$, AB = 1. If $\angle$BAD is acute and AD = a + b$\sqrt5$, then the value of a + b is equal to & (S) & 3 \\ \hline
 & & (T) & 4 \\ \hline
\end{tabular}}\par
\begin{qoptions}
\opt{A}{I $\to$ Q;II $\to$ P;III $\to$ R;IV $\to$ T}
\opt{B}{I $\to$ P;II $\to$ R;III $\to$ P;IV $\to$ R}
\opt{C}{I $\to$ T;II $\to$ R;III $\to$ Q;IV $\to$ S}
\opt{D}{I $\to$ S;II $\to$ P;III $\to$ R;IV $\to$ Q}
\end{qoptions}
\qn{9} Match the following List-I with List-II\par
{\footnotesize\setlength{\tabcolsep}{1.5pt}\noindent\begin{tabular}{|c|L{4.6cm}|c|c|}
\hline
 & \multicolumn{1}{>{\centering\arraybackslash}p{4.6cm}|}{List-I} & & List-II \\ \hline
(I) & If f(x) + f$''$(x) = $-$x g(x)f$\,'$(x) and g(x) $>$ 0 x $\in$ R then (f(x))$^2$ + [f$\,'$(x)]$^2$ has k as its maxima. Then k is & (P) & 0 \\ \hline
(II) & If f$\,^2$(x) + f(x) = g(x) and g(x) is always increasing then the minimum value that f(x) can attain so that f(x) also increasing is k. Then (2(k + 1)) is & (Q) & 1 \\ \hline
(III) & If g(x) = f(x)$\displaystyle\int_0^x$ f(t) dt is non-increasing function and f(x) is continuous in R then f(1) is & (R) & 2 \\ \hline
(IV) & If f(x) is differentiable function and f(x) = x$^2$ + $\displaystyle\int_0^x$ e$^{-t}$f(x $-$ t) dt then f$\,'$($-$1) = $\dfrac k3$ then k is & (S) & $-$1 \\ \hline
 & & (T) & 3 \\ \hline
\end{tabular}}\par
\begin{qoptions}
\opt{A}{I $\to$ P;II $\to$ Q;III $\to$ P;IV $\to$ R}
\opt{B}{I $\to$ Q;II $\to$ T;III $\to$ P;IV $\to$ R}
\opt{C}{I $\to$ S;II $\to$ Q;III $\to$ T;IV $\to$ S}
\opt{D}{I $\to$ R;II $\to$ T;III $\to$ P;IV $\to$ Q}
\end{qoptions}
\qn{10} Let H : $\dfrac{x^2}{a^2}-\dfrac{y^2}{b^2}$ = 1, where a,b $>$ 0 be a hyperbola in the xy-plane whose centre is C, one of the foci is F$_1$. Foot of perpendicular from F$_1$ upon line L : bx $-$ ay = 0 is P. Triangle CF$_1$P is isosceles triangle whose area is 8 sq. units.\par
Match List-I with List-II and select the correct answer using the code given below the list.\par
{\footnotesize\setlength{\tabcolsep}{1.5pt}\noindent\begin{tabular}{|c|L{3.6cm}|c|c|}
\hline
 & \multicolumn{1}{>{\centering\arraybackslash}p{3.6cm}|}{List-I} & & List-II \\ \hline
(I) & The length of transverse axis of H is & (P) & 4$\sqrt2$ \\ \hline
(II) & The eccentricity of H is & (Q) & 1 \\ \hline
(III) & The slope of L is & (R) & 4 \\ \hline
(IV) & The distance between directrices of H is & (S) & $\sqrt2$ \\ \hline
 & & (T) & 8 \\ \hline
\end{tabular}}\par
\begin{qoptions}
\opt{A}{I $\to$ S;II $\to$ T;III $\to$ Q;IV $\to$ P}
\opt{B}{I $\to$ T;II $\to$ S;III $\to$ Q;IV $\to$ P}
\opt{C}{I $\to$ T;II $\to$ S;III $\to$ Q;IV $\to$ R}
\opt{D}{I $\to$ T;II $\to$ R;III $\to$ Q;IV $\to$ P}
\end{qoptions}
\sect{SECTION-II}
\qn{1} Suppose f and g are differentiable function such that xg(f(x))f $\,'$(g(x))g$'$(x) = f(g(x))g$'$(f(x))f$'$(x) for all real x. Also f is non negative and g is positive. If $\displaystyle\int_0^a$ f(g(x)) dx = $\dfrac12-\dfrac{e^{-2a}}2$ for all reals a and g(f(0)) = 1 then the value of g(f(4)) is equal to e$^{-\lambda}$ where $\lambda\in$ N, then the value of $\dfrac\lambda{10}$ is\par
Note : (where `e' is napier's constant)
\qn{2} Let z satisfies z$\overline{\mathrm z}$ + ($-$4 + 5i)$\overline{\mathrm z}$ + ($-$4 $-$ 5i) z $-$ 40 = 0. If a = max.$|$z + 2 $-$ 3i$|$ and b = min$|$z + 2 $-$ 3i$|$, then value of $\left(\dfrac{a-b}8\right)$ is ?
\qn{3} Let f(1) + g(1) = 9e; f(x) = $-$x$^2$g$'$(x); g(x) = $-$x$^2$f$'$(x). If $\displaystyle\int_1^4\dfrac{f(x)+g(x)}{x^2}$dx = k$\left(e-e^{\frac14}\right)$, then the value of $\dfrac k4$ is\par
(Note : where e is napier's constant)
\qn{4} Graph of a function y = f(x) is shown as\par
\qfig{p1_math_s3_q04_1}
If g(x) = $|$f($|$x$|$)$|$, then find number of solution(s) of the equation\par
g(g(x)) = sgn(x$^2-$(k + 1)x + (k$^2$ + 1)), k $\in$ R.\par
[Note : sgn(y) denotes the signum function of y.]
\qn{5} Let D = $\begin{bmatrix}1&0&0\\0&2&0\\0&0&3\end{bmatrix}$ and P = $\begin{bmatrix}7&0&2\\0&1&0\\2&0&5\end{bmatrix}$. Consider A = P$^{-1}$DP then $\dfrac{\det.(A^2+A)}{10}$ is
\qn{6} Let A$_r$,r = 1,2,.....29 be arithmetic means between 303 and $-$57 where A$_r>$ A$_{r+1}$ $\forall$ r = 1,2,....28. If S be the sum of these means, then the value of $\left(\dfrac{S}{(A_{14}-12)|A_r|_{\min}}\right)$\par
[Note : $|$A$_r|_{\min}$. denotes the minimum value of $|$A$_r|$]
\qn{7} If $\left(x-2+\dfrac1x\right)^{30}$ = n$_0$x$^{30}$ + n$_1$x$^{29}$+...+n$_{29}$x + n$_{30}$ + n$_{31}$x$^{-1}$+...+n$_{60}$x$^{-30}$ and C = n$_0$ + n$_1$ + n$_2$+...+n$_{60}$.\par
Find the value of a if n$_{30}$ $-$ C = $\dbinom{60}{a}$ .\par
[Note : $\dbinom nr$ denotes $^n$C$_r$.]
\qn{8} Let E $-$ ABCD be a pyramid on square base ABCD where A is the origin and B and D are lying on positive x-axis and y-axis respectively. If E is (0,2,3) and $\overrightarrow{\mathrm{DE}}\cdot\left(\hat i+\hat j\right)=\vec 0$, then volume of the tetrahedron ABDE is
\raggedcolumns
\papersec{2}
\partsec{1}{PHYSICS}
\sect{SECTION-I(i)}
\qn{1} A horizontal tube of uniform cross-sectional area A is bent in the form of U as shown in figure. If the liquid of density $\rho$ enters and leaves the tube with velocity v, then the external force F required to hold the bend stationary is :-\par
\qfig{p2_phy_s1_q01_1}
\begin{qoptions}
\opt{A}{F = 0}
\opt{B}{$\rho$Av$^2$}
\opt{C}{2$\rho$Av$^2$}
\opt{D}{None of these}
\end{qoptions}
\qn{2} The activity of a radioactive sample goes down to about 6\% in a time of 2 hours. The half life of sample in minutes is about :-
\begin{qoptions}
\opt{A}{30}
\opt{B}{15}
\opt{C}{60}
\opt{D}{120}
\end{qoptions}
\qn{3} A planet of mass m is revolving in an elliptical orbit about the sun with an orbital period T. If M$_\mathrm{sun}$ be the mass of sun and A be the area of orbit, then the angular momentum of the planet is :-
\begin{qoptions}
\opt{A}{$\dfrac{\mathrm{mA}}{2\mathrm T}$}
\opt{B}{$\dfrac{\mathrm{mA}}{4\mathrm T}$}
\opt{C}{$\dfrac{2\mathrm{mA}}{\mathrm T}$}
\opt{D}{$\dfrac{4\mathrm{mA}}{\mathrm T}$}
\end{qoptions}
\qn{4} Figure shows a thin metal sheet in the plane y = 0, for which the current of constant density flows in the positive x-direction. It is in a constant homogeneous magnetic field of value B$_0\hat{\mathrm k}$. As a result of superposition of magnetic field in region y $>$ 0, the induction field B$_1\hat{\mathrm k}$ and in y $<$ 0 is B$_2\hat{\mathrm k}$ where B$_1 >$ B$_2$. Specify the correct statement :-\par
\qfig{p2_phy_s1_q04_1}
\begin{qoptions}
\opt{A}{B$_0 = \dfrac{(\mathrm B_1-\mathrm B_2)}2$}
\opt{B}{B$_0 = \dfrac{(\mathrm B_1+\mathrm B_2)}2$}
\opt{C}{B$_0$ = B$_1$ + B$_2$}
\opt{D}{B$_0$ = B$_1$ $-$ B$_2$}
\end{qoptions}
\sect{SECTION-I(ii)}
\qn{5} Consider a thermodynamics cycle in a PV diagram shown in the figure performed by one mole of a monoatomic gas. The temperature at A is T$_0$ and volumes are related as V$_\mathrm B$ = V$_\mathrm C$ = 2V$_\mathrm A$. Choose the correct options :-\par
\qfig{p2_phy_s2_q05_1}
\begin{qoptions}
\opt{A}{The maximum temperature during the cycle is 4T$_0$}
\opt{B}{Net work done by the gas during the cycle is 0.5 RT$_0$}
\opt{C}{The heat capacity of the process AB is 2R.}
\opt{D}{The efficiency of the cycle is 8.33\%}
\end{qoptions}
\qn{6} The plates of a parallel plate capacitor with no dielectric are connected to a voltage source. Now a dielectric of dielectric constant K is inserted to fill the whole space between the plates with voltage source remaining connected to the capacitor. Select the correct statement(s) :-
\begin{qoptions}
\opt{A}{The energy stored in the capacitor will become K-times the initial energy.}
\opt{B}{The electric field inside the capacitor will decrease K-times the initial field.}
\opt{C}{The force of attraction between the plates will become K$^2$-times the initial force.}
\opt{D}{The charge on the capacitor will become K-times the initial charge.}
\end{qoptions}
\qn{7} In a resonance tube experiment to determine speed of sound, air column in the pipe is made to resonate with a given fork of frequency 480 Hz. At room temperature T using a resonance tube a student finds two resonance position 16.0 cm and 51.0 cm, length of the air column for the first and the second resonance respectively. It is known that speed of the sound at 273 K is equal to 330 m/s. Then,
\begin{qoptions}
\opt{A}{percentage error in the measurement of speed of sound is $\dfrac47$\%}
\opt{B}{diameter of the tube is 5 cm.}
\opt{C}{temperature of the room T is nearly 283 K.}
\opt{D}{percentage error in the determination of T is $\dfrac87$\%}
\end{qoptions}
\sect{SECTION-II(i)}
\paraq{1}{2}
Two wave pulses are travelling in opposite direction with speed 1 ms$^{-1}$. Figure shows the shape of pulse at t = 0. Based on the above facts, answer the following questions.\par
\qfig{p2_phy_s3_para1_2_1}
\qn{1} The speed of particle (in m/s) at x = 2 cm and t = 0 is :-
\qn{2} The displacement of particle (in mm) at x = 8 cm and t = 0.06 sec is
\paraq{3}{4}
A small block of mass m = 1kg is placed over a wedge of mass M = 4kg as shown in figure. Mass m is released from rest. All surfaces are smooth and the origin O is fixed on the ground shown. Based on above information, answer the following questions.\par
\qfig{p2_phy_s3_para3_4_1}
\qn{3} Final velocity of the wedge is (in m/s)
\qn{4} The block strikes the x-axis at (in m)
\sect{SECTION-II(ii)}
\qn{5} In LCR circuit current resonant frequency is 600 Hz and half power points are at 650 and 550 Hz. The quality factor is
\qn{6} In a zero gravity region, a drop of a liquid of surface tension $\sigma$ assumes a cylindrical shape of diameter D between two parallel glass plates that are a distance `d' apart. If curved surface of the drop is at right angles to the plates as shown in the figure and force exerted by the drop on a plate is $\dfrac{\pi\sigma\mathrm D}{\mathrm n}$, then find `n'\par
\qfig{p2_phy_s4_q06_1}
\qn{7} A parallel beam of light is incident from air at an angle `$\alpha$' on the side PQ of a right angled triangular prism of refractive index n = $\sqrt2$. Light undergoes total internal reflection in the prism at the face PR when $\alpha$ has a minimum value of 45$^\circ$. The angle `$\theta$' of the prism is `n' times 3$^\circ$. Find `n'.\par
\qfig{p2_phy_s4_q07_1}
\qn{8} A wire having a length L and cross-sectional area A is suspended at one of its ends from a ceiling. Density and Young's modulus of material of the wire are $\rho$ and Y, respectively. Its strain energy due to its own weight is $\dfrac{\rho^2\mathrm g^2\mathrm{AL}^3}{\mathrm{pY}}$. Calculate the value of p.
\qn{9} In a hypothetical system a particle of mass m and charge $-$3q is moving around a very heavy particle having charge q. Assuming Bohr's model to be applicable to this system, the orbital velocity of mass m when it is nearest to heavy particle is $\dfrac{3\mathrm q^2}{\mathrm x\varepsilon_0\mathrm h}$. Find the value of x.
\qn{10} A cylinder of radius r = 0.1 m and mass M = 2 kg is placed such that it is in contact simultaneously with a vertical and a horizontal surface as shown in figure. The coefficient of static friction is $\mu$ = $\left(\dfrac13\right)$ for both the surfaces. Find the distance d, in cm, from the centre of the cylinder at which a force F = 40 N should be applied so that the cylinder just starts rotating in the anticlockwise direction. Take g = 10 ms$^{-2}$.\par
\qfig{p2_phy_s4_q10_1}
\raggedcolumns
\papersec{2}
\partsec{2}{CHEMISTRY}
\sect{SECTION-I(i)}
\qn{1} Consider the following unit cell of a metal:\par
\qfig{p2_chem_s1_q01_1}
Also consider the following shaded planes P and Q representing two 2D lattices.\par
\noindent P \qfigi{p2_chem_s1_q01_2}\quad Q \qfigi{p2_chem_s1_q01_3}\par
The ratio of packing fraction of P to Q is
\begin{qoptions}
\opt{A}{1.00}
\opt{B}{1.63}
\opt{C}{0.55}
\opt{D}{0.74}
\end{qoptions}
\qfullbegin
\qn{2} Vanillin (L), C$_8$H$_8$O$_3$, is isolated from vanilla beans. It give an intense colour with FeCl$_3$ and positive Tollen's test. It is not steam-distilled and does not react with HCl. It goes through following steps :\par
\qfig{p2_chem_s1_q02_0}
\qfullsplit
Which of the following is correct statement about various products?
\begin{qoptions}
\opt{A}{L is \qfigi{p2_chem_s1_q02_AL} and P is \qfigi{p2_chem_s1_q02_AP}}
\opt{B}{L is \qfigi{p2_chem_s1_q02_BL} and P is \qfigi{p2_chem_s1_q02_BP}}
\opt{C}{L is \qfigi{p2_chem_s1_q02_CL} and P is \qfigi{p2_chem_s1_q02_CP}}
\opt{D}{L is \qfigi{p2_chem_s1_q02_DL} and P is \qfigi{p2_chem_s1_q02_DP}}
\end{qoptions}
\qfullend
\qn{3} Predict the major product (P) in the following reaction :\par
\qfig{p2_chem_s1_q03_1}
\begin{qoptions}
\opt{A}{\qfigi{p2_chem_s1_q03_A}}
\opt{B}{\qfigi{p2_chem_s1_q03_B}}
\opt{C}{\qfigi{p2_chem_s1_q03_C}}
\opt{D}{\qfigi{p2_chem_s1_q03_D}}
\end{qoptions}
\qn{4} The correct molecular orbital diagram for O$_2$ molecule in the ground state is
\begin{qoptions}
\opt{A}{\qfigi{p2_chem_s1_q04_A}}
\opt{B}{\qfigi{p2_chem_s1_q04_B}}
\opt{C}{\qfigi{p2_chem_s1_q04_C}}
\opt{D}{\qfigi{p2_chem_s1_q04_D}}
\end{qoptions}
\sect{SECTION-I(ii)}
\qn{5} Consider one mole of a real gas A, with equation of state P(V $-$ b) = RT where `b' is a constant. Which of the following is/are correct statement(s) about A?\par
(P = pressure, V = volume, T = temperature in Kelvin, R = 8.314 J K mol$^{-1}$)
\begin{qoptions}
\opt{A}{Density of A is less than ideal gas at all pressure values.}
\opt{B}{Potential energy of A, due intermolecular forces, is either greater than (or) equal to that of ideal gas}
\opt{C}{If 5J of heat is given to A, and 2J work is done by the A then, its internal energy increases by 3J, which is true for ideal gas also}
\opt{D}{Compressibility factor (Z) of A is directly proportional to `P'.}
\end{qoptions}
\qn{6} When copper is reacted with dil. HNO$_3$ a gas is liberated. Which of the following are true for the gas?
\begin{qoptions}
\opt{A}{Has 11 valence e's, where extra electron occupies an antibonding $\pi$ *2p orbital}
\opt{B}{An acid anhydride}
\opt{C}{Paramagnetic}
\opt{D}{Neutral oxide}
\end{qoptions}
\qn{7} The complex(es), which can exhibit the type of isomerism shown by [Co(en)$_3$]Cl$_3$, is(are)\par
[en = H$_2$NCH$_2$CH$_2$NH$_2$]
\begin{qoptions}
\opt{A}{[Pt(en)(SCN)$_2$]}
\opt{B}{[Zn(gly)$_2$]}
\opt{C}{[Co(NH$_3$)$_4$Cl$_2$]$^+$}
\opt{D}{[Cr(en)$_2$(H$_2$O)(SO$_4$)]$^+$}
\end{qoptions}
\sect{SECTION-II(i)}
\paraq{1}{2}
Consider the following process of folding/unfolding of a tripeptide segment in a large protein. The free energy change ($\Delta$G) will depend on the interaction of the unfolded tripeptide with the solvent (water) and with the rest of the protein in the folded state (see below). Assume that the tripeptide is made up of one non-polar (hydrophobic, shaded) and two polar (hydrophilic; unshaded) residues. Use the following $\Delta$G data of formation of interaction between residue with water or protein :\par
i) a non-polar residue and the solvent (water) $\Delta$G = +8 kJ mol$^{-1}$\par
ii) a non-polar residue and the rest of the protein $\Delta$G = $-$4 kJ mol$^{-1}$\par
iii) a polar residue and the solvent (water) $\Delta$G = $-$16 kJ mol$^{-1}$\par
iv) a polar residue and the rest of the protein $\Delta$G = $-$14 kJ mol$^{-1}$\par
\qfig{p2_chem_s3_para1_2_1}
\qn{1} What is the magnitude of $\Delta$G for the folding of the tripeptide segment in KJ/mol?
\qn{2} What is the magnitude of $\Delta$G for the folding of the tripeptide segment if all the three residues are polar in KJ/mol?
\paraq{3}{4}
Given below is synthesis of nylon-66. Assume 100 \% yield in every step and none of the organic compounds involved are in excess.\par
\qfig{p2_chem_s3_para3_4_1}
\qn{3} The value of x is \underline{\hspace{1.5cm}}.
\qn{4} The value of y is \underline{\hspace{1.5cm}}.
\sect{SECTION-II(ii)}
\qn{5} A solution of 0.1 M CH$_3$COOH is placed between parallel electrodes of cross-section area 4 cm$^2$, separated by 2 cm. For this solution, resistance measured is 100 $\Omega$. The elevation in boiling point of the 0.1 M CH$_3$COOH solution is x $\times$ 10$^{-5}$. The value of x is (use following information)\par
K$_\mathrm b$ = 0.5K kg / mol; $\Lambda^\infty_\mathrm m$ (H$^+$) = 300 S cm$^2$ mole$^{-1}$; $\Lambda^\infty_\mathrm m$ (CH$_3$COO$^-$) = 100 S cm$^2$ mole$^{-1}$
\qn{6} H$_3$C $-$ C(=O) $-$ CH$_3$ $+$ I$_2$ $\xrightarrow{\text{NaOH}}$ Iodoform reaction\par
How many molecules of I$_2$ are required for 3 molecules of CH$_3$COCH$_3$?
\qn{7} Find total number of inner orbital diamagnetic complexes :\par
[Fe(CN)$_6$]$^{3-}$, [Co(H$_2$O)$_6$]$^{2+}$, [Cr(en)$_3$]$^{3+}$, [Fe(CN)$_6$]$^{4-}$, [Co(H$_2$O)$_6$]$^{3+}$, [Cr(CN)$_6$]$^{3-}$, [PtCl$_6$]$^{2-}$, [Ni(H$_2$O)$_6$]$^{2+}$, [Cr(CO)$_6$]
\qn{8} Consider following statements :\par
(1) CO$_2$(g) $+$ C(s) $\longrightarrow$ 2CO(g), it is endothermic reaction,\par
(2) In copper extraction from CuFeS$_2$, CaSiO$_3$ is main slag formed.\par
(3) Brass is an alloy consisting of Zn metal.\par
(4) Cu $+$ H$_2$SO$_4$(conc.) $\longrightarrow$ CuSO$_4$ $+$ H$_2$.\par
(5) In Hall-Heroult process, anode is corroded forming CO$_2$ gas only.\par
(6) Due to non-metallic impurites melting point of iron is raised.\par
(7) In extraction of copper from copper pyrite, roasting process is employed.\par
How many of the following are correct statements.
\qn{9} Consider the following precipitate/s\par
PbS , CuS , HgS , CdS , As$_2$S$_3$, As$_2$S$_5$, SnS, SnS$_2$, Sb$_2$S$_3$, Sb$_2$S$_5$\par
x = Total number of yellow ppt/s which soluble in yellow ammonium sulphide\par
y = Total number of black ppt/s which insoluble in hot \& dil. HNO$_3$\par
Find the value of $\dfrac{\mathrm x}{\mathrm y}$.
\qn{10} \qfigi{p2_chem_s4_q09_1}\par
Identify total number of stereoisomer's if double bond is present in cis configuration in given compound (X)
\raggedcolumns
\papersec{2}
\partsec{3}{MATHEMATICS}
\sect{SECTION-I(i)}
\qn{1} Consider f : [0, $\infty$) $\to$ [0, $\infty$) is bijective function and it satisfies f(x)$\cdot$e$^{\mathrm{f(x)}}$ = x, then the value of\par
$\displaystyle\lim_{x\to\infty}\dfrac{\mathrm f(x)}{\ln x}$ is
\begin{qoptions}
\opt{A}{0}
\opt{B}{1}
\opt{C}{e}
\opt{D}{Does not exist}
\end{qoptions}
\qn{2} Real number x, y, z satisfies\par
x + xy + xyz = 1\par
y + yz + xyz = 2\par
z + xz + xyz = 4\par
If the largest possible value of xyz is $\dfrac{\mathrm a+\sqrt{\mathrm b}}{\mathrm c}$, where b \& c are prime numbers, then (a + b + c) equals
\begin{qoptions}
\opt{A}{20}
\opt{B}{22}
\opt{C}{24}
\opt{D}{26}
\end{qoptions}
\qn{3} If y is implicit differentiable function of x such that y(x + y)$^2$ = x, then $\displaystyle\int\dfrac{dx}{x+3y}$ is equal to
\begin{qoptions}
\opt{A}{$\dfrac12\ln\left((x+y)^2+1\right)+\mathrm C$}
\opt{B}{$\ln\left((x+y)^2+1\right)+\mathrm C$}
\opt{C}{$\dfrac12\ln\left((x-y)^2+1\right)+\mathrm C$}
\opt{D}{$\ln\left((x-y)^2+1\right)+\mathrm C$}
\end{qoptions}
Where C is constant of integration
\qn{4} Let A = \{1,2, 3, 4, 5, 6, 7, 8\} be a set. Let f : A $\to$ A be a function such that f (x) $\neq$ x $\forall$ x $\in$ A, f (x) is even whenever x is even and f(x) is odd whenever x is odd. Then total number of onto functions satisfying these conditions will be
\begin{qoptions}
\opt{A}{36}
\opt{B}{49}
\opt{C}{64}
\opt{D}{81}
\end{qoptions}
\sect{SECTION-I(ii)}
\qn{5} Let f : R $\to$ R, g : R $\to$ R and h : R $\to$ R be differentiable function such that f(x) = x$^3$ + 5x + 1, g(f(x)) = x and h(g(g(x))) = x $\forall$ x $\in$ R then
\begin{qoptions}
\opt{A}{h(0) = 7}
\opt{B}{h$'$(1) = 1216}
\opt{C}{g$'$(1) = $\dfrac15$}
\opt{D}{h(g(2)) = 19}
\end{qoptions}
\qn{6} There are 20 points lying on a circle, distance between two consecutive points is same. Three points are selected at random. X represent no. of ways to select them. Y represent no. of ways to select them such that no two adjacent objects are selected and Z represent no. of ways to select such that chosen objects are not diametrically opposite then
\begin{qoptions}
\opt{A}{X = 1140}
\opt{B}{Y = 800}
\opt{C}{Z = 960}
\opt{D}{Y $\cap$ Z = 660}
\end{qoptions}
\qn{7} An ellipse having eccentricity $\dfrac34$ passes through the point P (5, 4). If the nearer focus is (1,1) and equation of tangent at P on the ellipse 4x + 3y $-$ 32 = 0, then
\begin{qoptions}
\opt{A}{Length of latus rectum of ellipse is 17.5}
\opt{B}{Length of chord through (1,1) parallel to the tangent at P is 17.5}
\opt{C}{Length of perpendicular from the focus other than (1, 1) on the tangent at P is 35}
\opt{D}{Locus of point of intersection of perpendicular tangents drawn to ellipse is (x + 1)$^2$ + (y + 8)$^2$ = 575}
\end{qoptions}
\sect{SECTION-II(i)}
\paraq{1}{2}
Let f(x) is strictly increasing and continuous function in [0, 1] and for x $\in$ [0, 1] g(x) = $\displaystyle\int_0^x$ f (t) dt, also\par
$\displaystyle\int_0^1$ (f (t) + $|$f (t) $|$) dt = 6. Local minimum value of g(x) is $-$1 occurs at x = a, where a $\in$ (0, 1), then
\qn{1} $\displaystyle\int_0^1$ g (x) f (x) dx is equal to
\qn{2} $\displaystyle\int_a^1$ f (x) dx is equal to
\paraq{3}{4}
A tetrahedron is formed by the planes P$_1$ : y + z = 0. P$_2$ : z + x = 0, P$_3$ : x + y = 0 and P$_4$ : x + y + z = 3.
\qn{3} The volume V of the tetrahedron is :
\qn{4} Square of distance of P$_4$ from the vertex opposite to it is
\sect{SECTION-II(ii)}
\qn{5} Consider f(x) = $\dfrac{x^2+4}{|x|}$ and g(x) = p$|$x$|$ + 2, p $>$ 0. If the exhaustive set of values of `p' for which f(x) $>$ g(x) $\forall$ x $\in$ R $-$ \{0\} is (a, b) then a + 4b is
\qn{6} A fair coin is tossed 10 times and the outcomes are listed. Let E$_i$ be the event that the i$^{\mathrm{th}}$ outcome is a head and E$_m$ be the event that the list contains exactly m heads. If E$_i$ and E$_m$ are independent, then `m' is equal to
\qn{7} A perpendicular is drawn from a point on the line $\dfrac{x-1}2=\dfrac{y+1}{-1}=\dfrac z1$ to the plane x + y + z = 3 such that the foot of the perpendicular Q also lies on the plane x $-$ y + z = 3, if the coordinate of Q is (a, b, c) then $|$a $-$ b $-$ c$|$ is
\qn{8} If $\vec{\mathrm a}_1$ = a$\hat{\mathrm i}$ + b$\hat{\mathrm j}$, $\vec{\mathrm a}_2$ = $\vec{\mathrm a}_1$ $-$ $\left(2\hat{\mathrm i}+\sqrt5\hat{\mathrm j}\right)$ and $\vec{\mathrm a}_3$ = $\vec{\mathrm a}_1$ + $\left(2\hat{\mathrm i}-\sqrt5\hat{\mathrm j}\right)$, then\par
$\begin{vmatrix}1+\vec a_1\cdot\vec a_1&1+\vec a_1\cdot\vec a_2&1+\vec a_1\cdot\vec a_3\\1+\vec a_2\cdot\vec a_1&1+\vec a_2\cdot\vec a_2&1+\vec a_2\cdot\vec a_3\\1+\vec a_3\cdot\vec a_1&1+\vec a_3\cdot\vec a_2&1+\vec a_3\cdot\vec a_3\end{vmatrix}$ is equal to
\qn{9} If $\displaystyle\sum_{r=0}^{3n}a_r(x-4)^r=\sum_{r=0}^{3n}A_r(x-5)^r$ and a$_k$ = 1 $\forall$ k $\geqslant$ 2n and $\displaystyle\sum_{r=0}^{3n}d_r(x-8)^r=\sum_{r=0}^{3n}B_r(x-9)^r$ and\par
$\displaystyle\sum_{r=0}^{3n}d_r(x-12)^r=\sum_{r=0}^{3n}D_r(x-13)^r$ and d$_k$ = 1 $\forall$ k $\geqslant$ 2n, then the value of $\dfrac{\mathrm A_{2n}+\mathrm D_{2n}}{\mathrm B_{2n}}$.
\qn{10} If $\theta$ is the angle between $\bar a$ and $\bar b$ such that $\sqrt3(\bar a\cdot\bar b)+\left|\bar b-(\bar a\cdot\bar b)\bar a\right|=1+\bar b\cdot\bar b$ where $|\bar a|$ = 1 then cosec$^2\theta$ is
\end{multicols}
\end{document}
